\exercisegroup

\begin{enumerate}
\def\labelenumi{\arabic{enumi})}
\item
  把 II, 第 83 页, 习题 8 的结果推广到空间 \(\mathcal{C}(\mathrm{T};\mathbf{C})\) 以及它们的有限维子空间。
\item
  证明：当 \(z\) 取遍单位圆盘 \(|z| \leqslant 1\) （\(\mathbf{C}\) 中的单位圆盘）时，由诸点 \((z, z^2, \dots, z^n)\) （空间 \(\mathbf{C}^n\) 中的点）生成的凸锥是整个空间 \(\mathbf{C}^n\) 。（注意，不可能存在不全为零的复数 \(c_k\) ( \(1 \leqslant k \leqslant n\) )，使得 \(\mathcal{R}\big(\sum_{k=1}^{n} c_k e^{ki\theta}\big) \geqslant 0\) 对 \(0 \leqslant \theta \leqslant 2\pi\) 成立；这里利用了这样的事实：\(\int_0^{2\pi}e^{ki\theta}d\theta = 0\) 对每个整数 \(k\neq 0.\) 成立。）
\item
  对 H（四元数的除环）上的拓扑向量空间，试给出与本节对应的定义和性质。
\end{enumerate}

